% !TEX root = misc-rerand.tex
% title: 杂记：再随机化

\section{零和掩码的作用}

Pre-sign 中，参与方 $i$ 将 Shamir 份额转换为衍生私钥的加法分片：
\begin{equation}
  x'_i:=\lambda(i,\mathbb{S})x_i+\zeta_i+\nabla x\cdot|\mathbb{S}|^{-1}.
\end{equation}
其中，$\zeta_i$ 是本次会话的零和掩码，满足 $\sum_{i\in\mathbb{S}}\zeta_i=0$。
它调整各方的分片，同时保持 $\sum_{i\in\mathbb{S}}x'_i=x+\nabla x$。
完整编排见「协议编排」的公式 \eqref{eq:orch:xi-derived}。

DKLS23 在第 3.1 节定义 $F_{\mathrm{Zero}}$，并在 Protocol 3.6 中用其输出再随机化私钥分片。
下文采用其“双方承诺并揭示种子，再按会话派生”的构造思路，以比特串异或合成共享种子。
原文见 \href{https://eprint.iacr.org/2023/765.pdf}{Threshold ECDSA in Three Rounds}。

\section{Keygen：共同生成共享种子}

对每个 $j\in\mathbb{K}\setminus\{i\}$，本方独立采样 256 比特贡献 $\epsilon^\circ_{i,j}$，
以及承诺盲化串 $\nu_{i,j}$，按公式 \eqref{eq:orch:seed-com} 生成承诺。
双方先交换承诺，收齐后再经双方之间的认证保密信道交换贡献与盲化串。
承诺绑定本次 Keygen 的会话号、双方编号及用途，揭示后必须验证。

双方验证通过后，按公式 \eqref{eq:orch:symmetric-blind} 合成
\begin{equation}
  \epsilon_{i,j}:=\epsilon^\circ_{i,j}\oplus\epsilon^\circ_{j,i}=\epsilon_{j,i}.
\end{equation}
每方对每个对端只保存一份合成后的共享种子，作为 Keyshare 的一部分。
先承诺再揭示，使一方无法在看到对方贡献后改变自己的贡献来指定合成结果。
若对方拒绝揭示、揭示长度不符或承诺验证失败，则中止本次 Keygen，不使用未完成的 Keyshare。

\section{Pre-sign：派生成对标量并反对称求和}

本次签名集合为 $\mathbb{S}\subseteq\mathbb{K}$，双方事先约定本次 Pre-sign 的会话号 $\mathrm{sid}$。
共享种子可以复用，但每次 Pre-sign 的会话号必须不同。
对每个 $j\in\mathbb{S}\setminus\{i\}$，双方各自计算
\begin{equation}
  \label{eq:misc-rerand:zeta-ij}\tag{zeta-ij}
  \begin{aligned}
    \zeta_{i,j}&:=\mathrm{HashToScalar}\bigl(
      \texttt{"zero\_mask"},\mathrm{sid},\mathbb{S}, \\
      &\hspace{5em}\min(i,j),\max(i,j),\epsilon_{i,j}\bigr)\in\mathbb{Z}_n.
  \end{aligned}
\end{equation}
这里 $\mathrm{HashToScalar}$ 视作输出到标量域的随机预言机，实现时可采用 XOF 与拒绝采样。
输入采用无歧义编码，其中 $\mathbb{S}$ 按编号递增排列。
由于双方使用相同的种子、会话信息和编号顺序，有 $\zeta_{i,j}=\zeta_{j,i}$。

本方再按编号大小取正负号，得到自己的零和掩码：
\begin{equation}
  \label{eq:misc-rerand:zeta}\tag{zeta}
  \zeta_i:=\sum_{\substack{j\in\mathbb{S}\\j<i}}\zeta_{i,j}
  -\sum_{\substack{j\in\mathbb{S}\\j>i}}\zeta_{i,j}.
\end{equation}
对每对参与方，编号较大的一方加上该标量，编号较小的一方减去同一个标量。
全员求和时各项两两抵消，因此 $\sum_i\zeta_i=0$。
派生只用已有共享种子和一致的公开输入，故本次 Pre-sign 无需为零和掩码增加通信。

\section{性质的适用范围}

零和关系是上述公式的代数性质，共享种子的生成则需要承诺、揭示与验证。
承诺约束贡献的选择时机，但不能阻止恶意方拒绝揭示，所以协议仍允许中止。
某一对参与方都知道彼此共享的种子；掩码对恶意联盟的隐藏性质还取决于联盟掌握哪些种子。
完整签名协议的多会话安全应依据原文的组合证明，不能仅由零和关系或每次使用不同会话号推出。
